% ============================================================================
% v0 TEMPLATE -- npj Computational Materials, "Article" type
% Built against the official Guide to Authors (npjcompumats-gta.pdf), checked
% 2026-08-28. Key constraints encoded below -- do NOT relax without checking
% the live guide first, since Nature Partner Journals updates these:
%
%   - Article type: unstructured abstract, max 250 words
%   - Main text: 5,000 words, EXCLUDING abstract/references/figures/tables
%   - Max 6 tables+figures combined
%   - References: max ~50, "as current as possible"
%   - Section order (MANDATORY, do not reorder):
%       Title page -> Abstract+keywords -> Introduction -> Results ->
%       Discussion -> Methods -> Acknowledgements -> Competing Interests ->
%       Contributions -> Funding -> References -> Figure legends -> Tables
%       -> Figures
%   - Title <= 150 characters. Running title <= 50 letters+spaces.
%   - LaTeX: any standard class (article.cls used here) + default Computer
%     Modern fonts. No visual formatting effort needed -- journal restyles
%     at production. Do NOT use non-standard fonts.
%   - Citations: NUMERICAL ONLY, superscript in text (not "Author (Year)").
%     "(ref. X)" instead of superscript only when a ref must sit next to a
%     number/equation/formula in running text.
%   - References must be pasted into this .tex file directly before
%     submission (per the guide: generate .bbl via BibTeX, copy its content
%     in, then delete \bibliography/\bibliographystyle). Left as natbib +
%     placeholder \bibitem entries below for now -- swap for the real list
%     when the reference set is finalized.
%   - Data/code availability statements go at the END of Methods, code
%     availability BEFORE data availability, per the guide.
%   - All pages and lines must be numbered at submission (line numbers via
%     lineno below; can be removed once accepted).
%
% This file is a STRUCTURAL v0: real numbers from the project's verified
% eval pipeline are pre-filled in the Results tables (see MASTER_SUMMARY.md
% in the project repo for full provenance of every number below); prose in
% Introduction/Discussion and all [TODO] markers still need writing.
% ============================================================================

\documentclass[11pt]{article}

\usepackage[utf8]{inputenc}
\usepackage[T1]{fontenc}
\usepackage{amsmath,amssymb}
\usepackage{graphicx}
\usepackage{booktabs}
\usepackage[super,comma,sort&compress]{natbib}   % gives superscript numeric cites, per house style
\usepackage[left]{lineno}                         % required at submission; strip after acceptance
\usepackage{setspace}
\doublespacing                                    % house style: double-spaced text + references
\usepackage[margin=1in]{geometry}

\linenumbers

\begin{document}

% ---------------------------------------------------------------------------
% TITLE PAGE
% ---------------------------------------------------------------------------

\title{Open-boundary reference data at interface scale: a coverage-driven
active-learning MACE potential for the Li-metal / LLZO solid-electrolyte
interface}

% Running title (<=50 letters and spaces):
% Open-boundary MLIP for Li/LLZO interfaces

\author{
Mehul Darak$^{1}$ \and
[Coauthor Name]$^{1,2}$ \and
[PI Name]$^{1,\ast}$
}

\date{}
\maketitle

\noindent
$^{1}$[Department, Institute, City, Country] \\
$^{2}$[Second affiliation if needed] \\
$^{\ast}$Corresponding author: [PI email] \\

\vspace{1em}

% ---------------------------------------------------------------------------
% ABSTRACT AND KEYWORDS
% ---------------------------------------------------------------------------
\section*{Abstract}
Machine-learned interatomic potentials (MLIPs) are now routinely used to
extrapolate from small periodic training cells to interface simulations of
several hundred to a few thousand atoms. At a Li-metal/garnet electrolyte
contact this extrapolation is doubly unsafe: the training cells are cheap
periodic approximations to a system whose defining feature is a free surface,
and the reference method itself --- plane-wave density functional theory (DFT)
--- cannot represent an isolated interface without a periodic-sandwich
construction, an image dipole, or a post-hoc correction. We remove both
approximations at once by generating every reference label with DFT-FE, a
real-space finite-element DFT code whose mixed boundary conditions permit
genuinely open (non-periodic) treatment along the surface normal, at the
420--1378 atom scale at which the MLIP is actually deployed. A MACE foundation
model is adapted to this data by low-rank (LoRA) fine-tuning over six
active-learning iterations (16$\rightarrow$108 labelled structures) using a
coverage-driven acquisition rule in the pretrained model's latent embedding
space. Component-wise force error on a held-out in-domain set of 33 structures
falls from 121.2 to 72.4~meV/\AA\ (RMSE) and 69.9 to 45.1~meV/\AA\ (MAE) from
the foundation model to the final iteration, with the last iteration's gain
significant under a paired test ($p=0.008$). Independently, pre-round posterior
variance in the same latent space predicts the realised DFT force error of the
newly labelled structures with $r=+0.85$, showing that the acquisition signal
measures real model deficiency rather than geometric novelty alone.
[TODO: closing sentence once the diffusion-coefficient validation with
independent MD seeds is complete.]

\vspace{0.5em}
\noindent\textbf{Keywords:} machine-learned interatomic potentials, active
learning, MACE, solid electrolyte interface, LLZO, Li-metal battery, DFT-FE,
open boundary conditions

% ===========================================================================
\section{Introduction}
% ===========================================================================

Garnet-structured Li$_7$La$_3$Zr$_2$O$_{12}$ (LLZO) is among the most
developed inorganic solid electrolytes for Li-metal batteries, combining a
room-temperature ionic conductivity in the mS\,cm$^{-1}$ range with nominal
stability against metallic lithium. Its practical performance is nonetheless
limited not by the bulk electrolyte but by the Li/LLZO contact, where
interfacial resistance, void formation on stripping, and dendrite penetration
along grain boundaries all originate. These are atomistic, interface-localised
phenomena, and resolving them requires dynamics on cells large enough to
contain a genuine interface together with the electrode and electrolyte
regions on either side of it --- in practice several hundred to a few thousand
atoms, well beyond the reach of direct \emph{ab initio} molecular dynamics.

The standard resolution is a machine-learned interatomic potential trained on
DFT labels. Recent foundation MLIPs, in particular the MACE family pretrained
on large bulk-periodic datasets such as OMat24, provide a strong starting
point and are routinely fine-tuned to a target chemistry with modest amounts
of new data. Two assumptions are usually inherited silently along with that
convenience, and both are violated at this interface.

The first concerns scale. MLIPs learn from local atomic environments, so a
model trained on small cells is assumed to transfer to large ones provided the
local environments match. At an interface they do not match: the environments
that matter --- undercoordinated surface oxygen, Li in the space-charge
region, the first oxide layer facing vacuum --- exist only in a cell that
actually contains a free surface. Training on small periodic approximants and
extrapolating to 500--1500-atom interface simulations therefore demands a
large and diverse dataset precisely because the relevant environments are
never directly sampled. In this work the training structures are the same size
as the production simulations, so the extrapolation in cell size is removed
rather than compensated for with data volume.

The second concerns boundary conditions, and is the more fundamental. A
plane-wave DFT code enforces three-dimensional periodicity. An asymmetric
Li$|$LLZO slab therefore requires either a symmetric Li$|$LLZO$|$Li sandwich,
which introduces a second interface that must be assumed equivalent to the
first, or a single slab with a large vacuum gap plus a dipole correction to
suppress the spurious interaction between periodic images across that gap.
Neither produces a free surface; both constrain the accessible slab thickness.
Prior first-principles studies of this interface have worked within these
constraints. DFT-FE, a real-space finite-element Kohn--Sham
code\citep{dftfe,Das2022}, removes them: it supports mixed boundary conditions
natively, so a slab can be periodic in-plane and genuinely open along the
surface normal, with zero Dirichlet conditions on the wavefunctions and the
electrostatic potential at the domain boundary. Because it is real-space and
GPU-parallel, it reaches the 450--1500 atom systems needed here. The reference
data underlying this potential is therefore open-boundary data at the scale of
the production simulation --- which, to our knowledge, no existing Li/LLZO
MLIP training set is.

This choice has a direct consequence for the learning problem. The foundation
model is pretrained on overwhelmingly bulk-periodic data, so adapting it to
open boundaries is a genuine transfer problem rather than a chemical
refinement. The error analysis bears this out: the outermost vacuum-facing
LLZO layer carries on average $2.05\times$ the frame-mean force error
(range $0.98$--$3.79\times$), i.e. the transfer gap is localised exactly where
the boundary condition differs from pretraining.

DFT-FE labels at this scale are expensive, which makes the acquisition
strategy consequential rather than cosmetic. We use a coverage-driven rule
operating on the pretrained model's own latent embedding space: candidate
structures harvested from MLIP-driven molecular dynamics are embedded,
farthest-point sampling seeded on the existing training set selects the
structures that most reduce the worst-case latent distance from any pool point
to its nearest labelled neighbour, and only those are sent for DFT-FE
evaluation. Recent independent work has established that pretrained-MLIP
latent representations are an effective general acquisition signal for
offline, pool-based active learning under DFT-cost scarcity, consistently
outperforming random acquisition\citep{vargaumbrich2026}; we do not repeat
that comparison here, and instead report a direct internal validation, showing
that the pre-round latent-space posterior variance of a candidate predicts the
force error the model will actually turn out to have on it once labelled.

This paper reports the resulting campaign and model. Section~\ref{sec:results}
presents the six-iteration active-learning trajectory, the force- and
energy-error trend across iterations on a held-out set, the spatial and
chemical localisation of the residual error, the validation of the acquisition
signal against realised DFT error, an explicit statement of the model's domain
of validity in vacancy concentration, and a training-length control confirming
the reported model is converged. Physical validation via Li-ion diffusion
coefficients is in progress and reported separately.

% ===========================================================================
\section{Results}
\label{sec:results}
% ===========================================================================

\subsection{Reference data: open boundaries at production scale}

All 108 training and 35 test structures are semi-periodic, with
$\mathrm{pbc}=[\mathrm{True},\mathrm{True},\mathrm{False}]$: periodic in $x$
and $y$, open in $z$. Each contains a single Li-metal/LLZO interface with two
genuine free surfaces, 420--1378 atoms, in-plane cell vectors of 11.3--26.4~\AA\
and a real-space domain extent along $z$ of 52--80~\AA. The $z$ extent is
finite-element domain padding, not image separation; there is no second
interface and no dipole correction.

Two internal checks bound the label quality. First, the electrostatic boundary
condition is immaterial at these geometries: recomputing representative
interfaces with Neumann (zero-flux) rather than zero-Dirichlet conditions on
the electrostatic potential changes the total energy by
$\mathcal{O}(10^{-6})$~Ha and the maximum force component by
$\sim5\times10^{-5}$~Ha/Bohr ($\sim$2.6~meV/\AA), some twenty times below the
converged model error. Second, an accidental duplicate pair within the test
set --- the same physical configuration labelled twice on differently sized
real-space domains --- differs by 0.06~meV/atom in energy and 0.12~meV/\AA\ in
force, i.e. DFT-FE label noise sits roughly fifty times below the model error
and is not limiting.

\subsection{Active-learning trajectory}

The campaign ran six iterations from a 16-structure initial fine-tuning set to
108 labelled structures, with the foundation model (never a previous
iteration's model) as the fine-tuning base at every step:
16 $\rightarrow$ 28 ($+12$) $\rightarrow$ 33 ($+5$) $\rightarrow$ 38 ($+5$)
$\rightarrow$ 47 ($+9$) $\rightarrow$ 77 ($+30$) $\rightarrow$ 108 ($+31$).
Every one of the 108 is an independent DFT-FE single-point calculation on a
420--1378 atom open-boundary interface.

The final round illustrates the acquisition rule concretely. The candidate
pool comprised 62{,}029 configurations, drawn from 50~ps MLIP-driven MD at
1100~K on 36 vacancy-doped interface structures together with a set of 1~ns
retest trajectories, after a minimum-pairwise-distance physicality filter that
removed none of them. Seeded on the 77 structures already labelled at that
point, the existing training set already covered 98.999\% of that pool within
the coverage threshold; only 31 further structures were required to reach full
coverage, and farthest-point sampling converged there exactly. All 31 came
from the vacancy-doped trajectories and none from the retest pool, and most
were early-trajectory frames (steps 0--1800 of 50{,}000), indicating that the
novel information was the defect geometry itself rather than deep phase-space
exploration --- the acquisition rule locating new chemistry rather than new
thermal noise. Figure~\ref{fig:latent} shows the pool, the seed, and the
selected structures in the latent space in which this decision was made.

The vac2pfu tier (2 vacancies per formula unit) was removed from the candidate
pool before selection, because 6 of its 12 structures were dynamically
unstable under MD, exceeding 3000~K within 81--248 steps. The model's domain
of validity in vacancy concentration was therefore fixed by construction at
$\leq$1 vacancy per formula unit; Section~\ref{sec:domain} quantifies what
happens outside it.

\subsection{Force- and energy-error trend across iterations}

Table~\ref{tab:force-trend} reports component-wise force RMSE and MAE for the
foundation model and every iteration on the in-domain held-out set of 33
structures. Both metrics fall essentially monotonically after it0, and the two
rank the iterations identically, so the MAE trend confirms rather than alters
the RMSE-based conclusion.

\begin{table}[htbp]
\centering
\caption{Force-error trend across active-learning iterations on the 33-structure
in-domain held-out set (29 clean test frames plus the 4 in-domain
vacancy holdouts; the 2 structures at 2 vacancies/f.u. are outside the trained
domain and are reported separately in Section~\ref{sec:domain}). RMSE and MAE
are component-wise over all atoms and all three Cartesian force components.
``max'' is the worst per-structure value in the set.}
\label{tab:force-trend}
\begin{tabular}{lcccc}
\toprule
Model & F RMSE mean & F RMSE max & F MAE mean & F MAE max \\
 & (meV/\AA) & (meV/\AA) & (meV/\AA) & (meV/\AA) \\
\midrule
universal (foundation) & 121.2 & 201.9 & 69.9  & 114.4 \\
it0                     & 188.1 & 330.4 & 100.8 & 168.8 \\
it1                     & 117.2 & 167.8 & 68.8  & 99.0  \\
it2                     & 94.6  & 163.7 & 57.7  & 100.5 \\
it3                     & 90.5  & 158.2 & 55.6  & 99.0  \\
it4                     & 83.7  & 128.9 & 51.5  & 81.9  \\
it5                     & 77.2  & 112.9 & 47.5  & 72.1  \\
\textbf{it6}            & \textbf{72.4} & \textbf{99.9} & \textbf{45.1} & \textbf{59.0} \\
\bottomrule
\end{tabular}
\end{table}

Three features of Table~\ref{tab:force-trend} deserve comment. First, it0 is
worse than the foundation model on this set. The first fine-tuning round used
16 structures, too few to adapt the model to open boundaries without degrading
the pretrained representation; the campaign recovers the foundation-model level
by it1 and improves on it thereafter. Reporting it0 rather than starting the
trend at it1 makes the cost of under-resourced fine-tuning explicit. Second,
the improvement is concentrated in the tail: from it0 to it6 the per-structure
maximum RMSE falls by 70\% and the across-structure standard deviation by 82\%,
compared with 62\% for the mean. The campaign is buying worst-case reliability,
which is what determines whether long MD runs remain stable, more than it is
buying average accuracy. Third, the final iteration's gain is real but modest:
a paired Wilcoxon signed-rank test on the 33 in-domain structures gives
$77.21\rightarrow72.37$~meV/\AA\ RMSE, a 6.3\% reduction, $p=0.008$, with it6
better on 24 of 33 structures.

Energy error behaves differently, and we report it as flat rather than
improving. On the same 33 structures it5 and it6 give 3.80 and 4.14~meV/atom
respectively; the difference is not significant ($p=0.50$ paired). Energy error
plateaued at approximately it2 and has not moved since. We note additionally
that the energy metric itself is near its numerical floor in this setting:
absolute total energies are of order $2.3\times10^5$~eV, for which the float32
unit in the last place is $1.6\times10^{-2}$~eV, so the difference
$|E_{\mathrm{pred}}-E_{\mathrm{ref}}|$ loses several digits to cancellation.
Repeating an identical evaluation of an identical model on identical data
returns 4.376, 4.280 and 4.369~meV/atom on successive runs, while the force
RMSE is reproduced bit-for-bit. Differences in reported energy error below
roughly 0.1~meV/atom are therefore not resolvable in float32 and should not be
interpreted, here or in comparable studies that do not state their working
precision.

An independent confirmation set of 7 snapshots drawn at 0/50/100/200/400/800/
1000~ps from a single 1~ns MD trajectory gives a different and less flattering
picture, which we state plainly: on this set it6 (80.0~meV/\AA\ RMSE,
49.5~meV/\AA\ MAE) is essentially tied with it4 (76.0, 46.9) rather than
clearly better, on both metrics. This set probes temporal rather than
configurational diversity --- all 7 frames come from one trajectory of one
interface --- so it is a weaker discriminator between late iterations than the
33-structure set, but the honest reading is that the it4$\rightarrow$it6 gain
does not reproduce on it.

\subsection{Where the residual error lives}

Pooling all atoms of the clean test frames, the residual force error is
strongly localised by both chemistry and geometry. The oxide framework carries
it: Zr and La together are 15.4\% of the atoms but hold 77.8\% of the
worst-1\% squared error, with mean component-wise error of 102.3 and
83.4~meV/\AA\ respectively, against 25.1~meV/\AA\ for Li, which is 47.8\% of
the atoms. Geometrically, the outermost vacuum-facing LLZO layer runs at
$2.05\times$ the frame-mean error on average, with a per-frame range of
$0.98$--$3.79\times$.

This is the signature the boundary-condition argument predicts. The
environments least represented in bulk-periodic pretraining are precisely
those at a free surface, and that is where the adapted model remains weakest.
It also indicates where further labelling effort would pay: a targeted round
biased towards surface-layer environments rather than another broad sweep of
the configuration pool.

\subsection{The acquisition signal predicts realised DFT error}
\label{sec:uq}

A recurring weakness of diversity-based acquisition is that latent-space
distance may measure only geometric novelty, not model deficiency. We tested
this directly using data already generated by the campaign, with no additional
DFT.

For each candidate we computed a Gaussian-process posterior variance in the
fine-tuned model's latent feature space, conditioned on the 77 structures
labelled before the final round. Compared against the force and energy error
the previous model was subsequently measured to have on the 35 structures that
round labelled, the pre-round posterior variance predicts realised error with
Pearson $r=+0.847$ for force and $+0.777$ for energy ($\lambda=0.1$;
$r=+0.823/+0.785$ at $\lambda=0.01$; all $p\leq2.3\times10^{-8}$). The
correlation is not an artifact of contrasting selected against random
structures --- restricting to the 31 farthest-point-sampled structures alone
gives $r=+0.839$ force, $+0.784$ energy --- nor of any single leverage point
(leave-one-out range $[+0.817,+0.867]$; bootstrap 95\% CI $[+0.69,+0.93]$), nor
of structure size ($r=-0.03$ between atom count and error; the partial
correlation controlling for atom count is unchanged at $+0.850$).

The structures the acquisition rule chose were also high-uncertainty by this
independent criterion: their mean posterior variance was 0.233 against 0.064
for the pool ($\lambda=0.01$, Mann--Whitney $p=4.4\times10^{-12}$), and the
selected set exceeded all 1000 randomly drawn 31-structure subsets. This
agreement is non-trivial, because the two criteria rank the pool almost
independently: Spearman correlation between posterior variance and the
farthest-point-sampling criterion over the full 62{,}029-frame pool is only
$+0.10$ to $+0.16$, and the top 31 by posterior variance overlap the selected
31 by just 10--11 structures. The two rules agree in the extreme tail while
disagreeing over the bulk of the pool.

We deliberately do not claim that posterior variance is a better acquisition
criterion than farthest-point sampling. A head-to-head comparison on the same
35 ground-truth structures gives $r=+0.847$ versus $+0.724$ for force, which is
not significant under Williams' test for dependent correlations ($p=0.066$),
and most of the apparent advantage is attributable to feature standardisation
rather than to the criterion: running posterior variance in the
farthest-point-sampling preprocessing space drops it to $+0.683$, statistically
indistinguishable from plain minimum-distance ($p=0.65$). The one clean
methodological finding is that per-feature standardisation improves
error prediction over the L2-normalise-then-PCA preprocessing used here
($+0.611\rightarrow+0.716$ for energy, $p=0.010$). Moreover, the comparison
arena is biased: all 35 ground-truth points are the selected structures plus 4
random frames, so it can only measure how well each criterion ranks difficulty
\emph{within the region already chosen}, never what the alternative would have
selected instead. The defensible conclusion is that latent geometry genuinely
tracks model deficiency, and that the two signals are complementary --- a
per-structure informativeness score paired with an explicit batch-diversity
constraint.

\subsection{Domain of validity}
\label{sec:domain}

The model's training set contains no structure above 1 vacancy per formula
unit: of the 108 training structures, 12 are at 0.5 and 19 at 1 vacancy per
formula unit, and 77 are undoped. This was a construction decision, not a
post-hoc exclusion --- the 2-vacancy tier was removed from the candidate pool
at the MD stage for dynamical instability, before any selection or labelling
took place.

Two 2-vacancy structures nonetheless entered the held-out set, and we report
them as an explicit extrapolation probe rather than deleting them. On these,
force RMSE is 252.6~meV/\AA\ for it6 (266.3 for it5), against 72.4~meV/\AA\
in-domain --- a 3.5-fold degradation. The potential should not be used above
1 vacancy per formula unit, and this is the measurement that states the
boundary.

\subsection{Training-length control}

To confirm that the reported model is limited by data rather than optimisation,
the final iteration was retrained at 2000 epochs (stochastic weight averaging
from 1200) and 2500 epochs (averaging from 1600), against the production
setting of 1500 epochs (averaging from 1000), all other hyperparameters and
both data files held identical. Force error is unchanged: 72.04, 72.00 and
71.84~meV/\AA\ RMSE and 44.89, 44.87 and 44.78~meV/\AA\ MAE respectively.

The cause is visible in the training logs: both longer runs selected the same
pre-averaging checkpoint, at epoch 453, and their best post-averaging
checkpoints fell at epochs 1226 and 1896, after which validation loss plateaued
and then drifted upward. Additional epochs neither help nor harm. The reported
model is converged with respect to training length, and the remaining error is
a property of the data and the architecture.

\subsection{Diffusion-coefficient validation}

[TODO --- pending. Preliminary single-seed estimates from it6-driven MD at
1100~K give in-plane diffusion coefficients of $\sim$9--11$\times10^{-7}$~cm$^2$/s
in the LLZO region and $\sim$4.1$\times10^{-4}$~cm$^2$/s in the Li-metal region,
against literature-extrapolated values of $\sim$2.3$\times10^{-6}$ and
$\sim$6$\times10^{-4}$--1.2$\times10^{-3}$~cm$^2$/s respectively --- agreement
within a factor of 2--3 for both regions. These are single-trajectory numbers
with no error bars and are NOT reported as a result in this draft. Multiple
independent MD seeds per condition are in progress; this subsection, the
corresponding figure, and the closing sentence of the Abstract are to be
written once those complete.]

% ===========================================================================
\section{Discussion}
% ===========================================================================

The case for this potential rests less on its aggregate error than on what the
error was measured against. A force RMSE of 72~meV/\AA\ is unremarkable by the
standards of bulk MLIP benchmarks; what is unusual is that it is measured
against open-boundary DFT labels on 420--1378-atom cells containing a single
interface with two real free surfaces, which is the same object the model is
deployed on. The usual reporting convention --- train on small periodic
approximants, validate on small periodic approximants, deploy on large
interfaces --- leaves the extrapolation that actually matters unmeasured. Here
there is no such extrapolation to make, and the error analysis shows why it
would have been dangerous: the residual error concentrates by a factor of two
on exactly the vacuum-facing layer that a periodic training cell does not
contain.

This reframes what DFT-FE contributes. It is not simply a larger-capacity DFT
code; its mixed boundary conditions make a class of reference data available
that plane-wave codes cannot produce at all, and that class is the one an
interface MLIP needs. The cost is real --- each label is a single-point
calculation on 64 GPU nodes --- which is precisely why the acquisition strategy
carries weight rather than being a convenience.

On that strategy, the evidence here is stronger than a coverage argument alone.
The pre-round latent posterior variance predicts the force error the previous
model turns out to have on newly labelled structures at $r=+0.85$, robust to
leave-one-out, bootstrap, and the obvious confounds. This says the latent
geometry is tracking genuine model deficiency and not merely geometric
novelty --- the standard objection to diversity-based acquisition --- and it is
measured on real DFT labels rather than on a synthetic proxy. It is also
convergent evidence, in that the selected structures were extreme by a
criterion that ranks the pool nearly independently ($\rho\approx0.1$--$0.16$).
We resist the stronger claim that posterior variance should replace
farthest-point sampling: the comparison is not significant for forces, most of
its apparent margin is preprocessing rather than criterion, and the evaluation
arena consists largely of structures the incumbent method chose, which cannot
reveal what the challenger would have selected. The actionable finding for a
future round is narrower and more robust --- per-feature standardisation of the
latent space improves error prediction independently of which criterion is
applied on top of it.

The campaign's own trajectory argues that it should stop rather than continue
broadly. By the final round, the 77 already-labelled structures covered 98.999\%
of a 62{,}029-frame candidate pool, and 31 additions exhausted the remainder.
The pool is not merely expensive to exploit further, it is close to
uninformative; a seventh broad iteration would buy little. What the error
localisation does suggest is a targeted round --- surface-layer environments,
where the residual is concentrated --- rather than another undirected sweep.

Three limitations bound the claims. First, the energy channel has been flat
since approximately it2 and did not improve in the final round; part of this is
genuine, and part is that the reported energy metric sits near its float32
numerical floor at these absolute energies, which we quantify rather than
assume. Second, the force improvement is statistically established on the clean
structures ($p=0.008$ overall, $p=0.007$ on the clean 29 alone) but not on the
4 in-domain vacancy holdouts, where $n=4$ supports no claim in either
direction; that adding vacancy training data improved performance on pristine
interfaces is reported as observed, not as the designed outcome. Third, the
held-out set is largely in-distribution --- 21 of 29 clean frames are siblings
of an earlier iteration's selection pool --- so it is a legitimate held-out
test but not an out-of-distribution one, and the 7-snapshot temporal
confirmation set does not reproduce the it4$\rightarrow$it6 gain. Physical
validation through diffusion coefficients with proper statistics over
independent seeds remains the main outstanding item and is deliberately not
claimed here.

Finally, the model carries a stated domain of validity in vacancy
concentration, fixed at $\leq$1 vacancy per formula unit when the 2-vacancy
tier proved dynamically unstable at the sampling stage. The 3.5-fold error
increase measured outside it is reported as the boundary rather than suppressed.

% ===========================================================================
\section{Methods}
% ===========================================================================

\subsection{System and boundary conditions}
All structures are semi-periodic: periodic in $x$ and $y$, open
(non-periodic) in $z$. This is possible because reference labels come from
DFT-FE (real-space DFT) rather than a plane-wave code, giving a single
interface with two genuine free surfaces and no periodic-image interaction
along the surface normal --- avoiding the periodic-sandwich workaround
(second interface, image dipole or correction, constrained slab thickness)
standard in plane-wave DFT treatments of this
interface\citep{dftfe,Das2022}. The condition is enforced in code at every
stage of the pipeline and verified on all 108 training and 35 test frames.
The finite $z$ extent of each cell (52--80~\AA, with 12.7--36~\AA\ beyond the
atoms) is real-space domain padding for the finite-element solver, not image
separation.

Interfaces were constructed by lattice-matching Li-metal and LLZO surface
slabs, sampling the interfacial separation, and selecting the optimal gap by
energy; the LLZO reference state throughout is the ordered tetragonal
$I4_1/acd$ polymorph\citep{Murugan2007}. Vacancy-doped variants were generated
at 0.5, 1 and 2 vacancies per formula unit. [TODO: cite or restate the
lattice-matching and gap-optimisation procedure in enough detail to be
reproducible, or point to a deposited script.]

\subsection{Reference labels: DFT-FE}
All reference energies and forces are single-point Kohn--Sham DFT calculations
performed with DFT-FE\citep{dftfe,Das2022} using the PBE generalized-gradient
exchange--correlation functional\citep{pbe1996} and SPMS norm-conserving
pseudopotentials\citep{spms} for Li, La, Zr and O, applied consistently across
the entire dataset. The Kohn--Sham equations were discretized on a high-order
finite-element basis of polynomial order 7, with an automatically generated
mesh of size parameter 1.2~\AA\ refined within 8.0~\AA\ of each atomic core.
Self-consistency used Anderson mixing with Fermi--Dirac smearing at an
electronic temperature of 500~K and a convergence tolerance of
$5\times10^{-5}$~Ha. Boundary conditions were periodic in $x$ and $y$ and zero
Dirichlet in $z$ for both the wavefunctions and the electrostatic potential.
No structural relaxation was performed; the objective is accurate force labels
at fixed geometry.

The sensitivity of the labels to the electrostatic boundary condition was
checked by recomputing representative interfaces with Neumann (zero-flux)
conditions on the electrostatic potential in place of zero Dirichlet: total
energies changed by $\mathcal{O}(10^{-6})$~Ha and the maximum force component
by $\sim5\times10^{-5}$~Ha/Bohr ($\sim$2.6~meV/\AA).

Calculations were run on the OLCF Frontier system, distributed across 64 nodes
(4 AMD MI250X GPUs per node, one MPI task per GPU) in ground-state mode with
GPU acceleration and memory optimisation enabled. All production labels used
identical settings.

\subsection{Model architecture and fine-tuning}
The potential is a MACE equivariant message-passing
network\citep{batatia2023mace}, adapted from the \texttt{mace-omat-0-medium}
foundation model pretrained on OMat24\citep{omat24}. Every iteration
fine-tunes the foundation model directly, never a previous iteration's
weights, so that iterations differ only in training data.

Adaptation uses low-rank adaptation (LoRA)\citep{LoRA2021} to limit
catastrophic forgetting of the pretrained bulk chemistry on a small interface
dataset. For a pretrained weight matrix $W\in\mathbb{R}^{d\times d}$ the update
is constrained to a rank-$r$ subspace,
\begin{equation}
W' = W + \frac{\alpha}{r}BA,
\qquad A\in\mathbb{R}^{r\times d},\; B\in\mathbb{R}^{d\times r},\; r\ll d,
\end{equation}
so that the layer evaluates
$h_{\mathrm{out}} = Wh_{\mathrm{in}} + \tfrac{\alpha}{r}B(Ah_{\mathrm{in}})$.
Adapters are inserted in the equivariant \texttt{o3.Linear} and dense layers
using compatible irreducible representations, so the low-rank update preserves
$O(3)$ equivariance. All pretrained parameters are frozen
($\nabla_{\theta_{\mathrm{base}}}L = 0$) and only $(A,B)$ are optimised,
reducing trainable parameters per layer from $\mathcal{O}(d^2)$ to
$\mathcal{O}(2rd)$, under 1\% of the model. Initialisation is
$A\sim\mathcal{N}(0,10^{-3})$, $B=0$, so that $\Delta W = 0$ at the start and
the model begins from exactly the pretrained state. For deployment the update
is merged, $W_{\mathrm{final}} = W + BA$, giving a standard MACE model with no
inference overhead.

Training minimises
\begin{equation}
L = w_E\sum_i\left(E_i-\hat{E}_i\right)^2
  + w_F\sum_i\left\|F_i-\hat{F}_i\right\|^2 ,
\end{equation}
with $w_E = 1.0$ and $w_F = 5.0$. Hyperparameters, identical across all
iterations: LoRA rank $r=5$, $\alpha=1$; batch size 8; learning rate
$5\times10^{-3}$ with \texttt{ReduceLROnPlateau} (factor 0.8, patience 50);
1500 epochs with stochastic weight averaging from epoch 1000 at learning rate
$10^{-4}$ and unchanged loss weights; seed 42; float32 working precision.
Isolated-atom reference energies (eV) were
$E_0(\mathrm{Li})=-190.7590256408$, $E_0(\mathrm{O})=-442.9888796243$,
$E_0(\mathrm{Zr})=-1380.1817128081$, $E_0(\mathrm{La})=-958.0774205521$.
Checkpoint selection is by best validation loss.

Working precision is float32 throughout, and this is a deliberate choice
rather than the library default (MACE defaults to float64). It has one
reporting consequence, quantified in Results: at absolute total energies of
order $2.3\times10^5$~eV the float32 unit in the last place is
$1.6\times10^{-2}$~eV, so the energy-error metric carries run-to-run scatter of
order 0.1~meV/atom while force errors reproduce exactly.

\subsection{Molecular dynamics sampling}
Candidate configurations were generated by MLIP-driven MD in ASE using the
current iteration's model. Simulations used a Langevin thermostat at 1100~K
with friction $10^{-3}$~fs$^{-1}$, a 1~fs timestep, and 50~ps (50{,}000 steps)
per structure, with configurations recorded every 100 steps. Boundary
conditions match the reference data exactly ($\mathrm{pbc}=[\mathrm{T},
\mathrm{T},\mathrm{F}]$), and only the $x,y$ coordinates are wrapped.

To prevent the slab from drifting and to confine sampling to the interfacially
active region, an active region was defined geometrically: the LLZO framework
was identified from La, Zr and O atoms, Li metal as the Li above it, and oxide
atoms deeper than $\max(0.85\,t,\,6~\text{\AA})$ below the LLZO top surface,
where $t$ is the slab thickness, were held fixed. This stabilises the far-bulk
region while leaving the interface free to evolve.

Generated frames passed a physicality filter (PBC-aware minimum pairwise
distance $>1.2$~\AA) before entering the candidate pool; in the final round
0 of 62{,}029 frames were rejected.

\subsection{Latent embeddings and coverage-driven acquisition}
For each structure, a forward pass with force evaluation enabled yields for
each atom $i$ the rotationally invariant ($l=0$) scalar components of the final
message-passing layer, $\mathbf{z}_i\in\mathbb{R}^{256}$. A structure-level
embedding is obtained by mean-pooling,
\begin{equation}
\mathbf{Z}_S = \frac{1}{N_S}\sum_{i=1}^{N_S}\mathbf{z}_i ,
\end{equation}
where $N_S$ is the atom count. Embeddings are L2-normalised and reduced by PCA
retaining at least 95\% of variance (minimum 32 components).

Selection is farthest-point sampling seeded on the existing labelled set
$\mathcal{S}_0 = \mathcal{T}^{(n)}$, so that new picks complement rather than
duplicate what is already labelled. For each pool candidate $k$,
\begin{equation}
\delta_k = \min_{\mathbf{Z}_s\in\mathcal{S}} \left\|\mathbf{Z}_k-\mathbf{Z}_s\right\|_2 ,
\end{equation}
and the candidate with the largest $\delta_k$ is added to $\mathcal{S}$,
repeatedly. The quantity being minimised is the coverage radius
\begin{equation}
\rho(\mathcal{S},\mathcal{P}) = \max_{k\in\mathcal{P}} \delta_k ,
\end{equation}
the worst-case distance from any pool point to its nearest selected neighbour,
which greedy farthest-point sampling reduces as fast as possible at each step.
Selection stops when the percentile coverage
\begin{equation}
C(d) = \frac{1}{|\mathcal{P}|}\sum_{k\in\mathcal{P}}\mathbf{1}\!\left[\delta_k\leq d\right]
\end{equation}
reaches the configured target (1.0 in the final round) at the threshold
distance $d$. This is a stopping rule driven by the pool's own geometry rather
than a fixed per-round budget, which is why round sizes vary from 5 to 31
structures.

For the posterior-variance analysis of Section~\ref{sec:uq}, a Gaussian-process
posterior variance was computed in the same feature space, conditioned on the
pre-round training set, at regularisation $\lambda\in\{0.01,0.1\}$;
significance of differences between dependent correlations used Williams' test.

\subsection{Test-set construction}
The held-out set comprises 35 structures. Twenty-nine come from an earlier
evaluation set of 32 frames after removing three contaminated entries,
identified not by filename but by model prediction: two test frames were
confirmed to be the same physical configuration as training frames recomputed
on a larger real-space domain (predicted energies agreeing to
$6.5\times10^{-6}$ and $1.8\times10^{-6}$~meV/atom), and one further pair was
an internal duplicate. The remaining 6 are a holdout set constructed
independently of the selection round: 4 random frames drawn from the
physicality-filtered 0.5/1-vacancy pool excluding all selected structures, and
2 taken at their pristine $t=0$ geometry from the 2-vacancy tier.

Because the model's training domain excludes 2 vacancies per formula unit
entirely (Section~\ref{sec:domain}), the primary reported metric is the
33-structure in-domain set (29 clean + 4 in-domain vacancy holdouts), and the
2 out-of-domain frames are reported separately as an extrapolation probe.
Twenty-one of the 29 clean frames are siblings of an earlier iteration's
selection pool, so this is a held-out but largely in-distribution set; it is
not an out-of-distribution benchmark and is not presented as one.

Force error is component-wise throughout:
$\mathrm{RMSE}=\sqrt{\langle(F_{\mathrm{pred}}-F_{\mathrm{ref}})^2\rangle}$
averaged over all atoms and all three Cartesian components, in meV/\AA;
multiply by $\sqrt{3}$ for the RMS of per-atom error vectors. Energy error is
$|E_{\mathrm{pred}}-E_{\mathrm{ref}}|/N$ in meV/atom. Paired model comparisons
use the Wilcoxon signed-rank test over structures.

\subsection{Molecular dynamics analysis}
Diffusion coefficients are extracted from in-plane ($x,y$) mean-squared
displacement with multiple time origins and an Einstein-relation fit,
$D_{2\mathrm{D}} = \mathrm{slope}/4$. Because trajectory frames are stored
wrapped in $x$ and $y$, coordinates are unwrapped before any displacement is
computed, by minimum-imaging successive per-frame displacements in fractional
coordinates and accumulating:
$\Delta\mathbf{f} \leftarrow \Delta\mathbf{f} - \mathrm{round}(\Delta\mathbf{f})$,
then converting back to Cartesian. The fractional formulation is required
because not all in-plane cells are orthogonal. Unwrapping is unambiguous at the
100~fs snapshot interval, where the RMS Li displacement per frame
($\sim$0.2~\AA) is far below half the smallest in-plane cell vector
(5.6~\AA); without it the MSD saturates near $L^2/6\approx21$~\AA$^2$ per
dimension in the smallest cells, which is reached within tens of picoseconds
at 1100~K.

The $z$ coordinate is excluded from the diffusion fit, being confined by the
slab's finite thickness rather than a freely diffusing direction, and is used
only to assign atoms to the LLZO or Li-metal region. That assignment is made
relative to the measured interface plane, defined per frame as the maximum $z$
of the La/Zr/O framework, rather than at a fixed absolute height, so that
region definitions remain valid as the slab breathes and across structures with
different cell extents. Atoms crossing between regions within a lag window are
excluded from that window for both regions. Atoms frozen by the active-region
constraint are excluded from the mobile population; no centre-of-mass drift
correction is applied, as the frozen substrate already fixes the reference
frame and subtracting an all-atom centre of mass would remove part of the
genuine Li flux.

[TODO: fitting window, number of independent seeds, temperature set, and
Arrhenius extrapolation for the literature comparison --- to be completed with
the multi-seed runs.]

\subsection*{Code availability}
[TODO: repository / archive link for training, evaluation, and MD/analysis
scripts.]

\subsection*{Data availability}
[TODO: repository / archive link and persistent identifier (DOI) for
DFT-FE reference labels, training/test structure files, and trained model
checkpoints.]

% ===========================================================================
\section*{Acknowledgements}
% ===========================================================================
[TODO: funding sources and named acknowledgements.] Model training, inference
and molecular dynamics were carried out on the KIAC Mahamathi cluster (NVIDIA
H200 and A5000 GPUs); DFT-FE single-point calculations were performed on the
OLCF Frontier supercomputer.

% ===========================================================================
\section*{Competing Interests}
% ===========================================================================
The authors declare no competing interests.

% ===========================================================================
\section*{Contributions}
% ===========================================================================
[Author 1] designed and ran the active-learning campaign, model training,
and evaluation, and wrote the manuscript. [Author 2] [TODO]. [PI] supervised
the project and is the guarantor of this work.

% ===========================================================================
\section*{Funding}
% ===========================================================================
[TODO: grant/funding-body statement, or "Not applicable" if none.]

% ===========================================================================
% REFERENCES
% ===========================================================================
\begin{thebibliography}{99}

\bibitem{batatia2023mace}
Batatia, I. \emph{et al.} MACE: Higher order equivariant message passing
neural networks for fast and accurate force fields. \emph{arXiv preprint}
arXiv:2206.07697 (2023).

\bibitem{omat24}
[TODO: OMat24 / mace-omat-0-medium foundation-model citation.]

\bibitem{dftfe}
[TODO: DFT-FE methodology citation (Motamarri et al., Comput. Phys. Commun.).]

\bibitem{Das2022}
Das, S. \emph{et al.} [TODO: complete DFT-FE GPU / large-scale accuracy
citation establishing agreement with plane-wave norm-conserving results.]

\bibitem{pbe1996}
Perdew, J. P., Burke, K. \& Ernzerhof, M. Generalized gradient approximation
made simple. \emph{Phys. Rev. Lett.} \textbf{77}, 3865--3868 (1996).

\bibitem{spms}
[TODO: SPMS norm-conserving pseudopotential citation
(\texttt{github.com/SPARC-X/SPMS-psps}).]

\bibitem{LoRA2021}
Hu, E. J. \emph{et al.} LoRA: Low-rank adaptation of large language models.
\emph{arXiv preprint} arXiv:2106.09685 (2021).

\bibitem{Murugan2007}
Murugan, R., Thangadurai, V. \& Weppner, W. Fast lithium ion conduction in
garnet-type Li$_7$La$_3$Zr$_2$O$_{12}$. \emph{Angew. Chem. Int. Ed.}
\textbf{46}, 7778--7781 (2007).

\bibitem{burov2026}
Burov, D., Abakumov, A. \& Aksyonov, D. Machine learning assisted molecular
dynamics of charge-transfer mechanisms at Li/Ga-doped LLZO interfaces.
\emph{arXiv preprint} arXiv:2606.07772 (2026).

\bibitem{vargaumbrich2026}
Varga-Umbrich, E., Surana, S., Duckworth, P., Tilly, J., Peltre, O. \&
Weller-Davies, Z. Pretrained model representations as acquisition signals
for active learning of MLIPs. In \emph{Proceedings of the 43rd International
Conference on Machine Learning (ICML 2026)}. Preprint: arXiv:2605.03964.

% [TODO: remaining references -- prior Li/LLZO interface DFT literature
% (Iwasaki et al., Burov et al. 2024), interfacial-resistance and dendrite
% experimental context, ASE, foundation-MLIP softening. Keep <=50 total.]

\end{thebibliography}

% ===========================================================================
\section*{Figure Legends}
% ===========================================================================

\paragraph{Figure 1.} System and boundary conditions. A single Li-metal/LLZO
interface with two genuine free surfaces, periodic in $x,y$ and open in $z$,
contrasted with the periodic-sandwich and dipole-corrected slab constructions
required by plane-wave DFT. [TODO: build figure.]

\paragraph{Figure 2.} \label{fig:latent} Latent-space view of the final
acquisition round. PCA on L2-normalized 256-dimensional MACE mean-pooled
embeddings (PC1/PC2/PC3 explain 70.9\%/21.8\%/2.9\% of variance), showing the
62{,}029-frame candidate pool, the 77-structure training seed, and the 31
selected structures.

\paragraph{Figure 3.} Force RMSE and MAE across active-learning iterations
(universal, it0--it6) on the 33-structure in-domain held-out set, corresponding
to Table~\ref{tab:force-trend}.

\paragraph{Figure 4.} Residual force error versus distance from the free
surface, pooled over the clean held-out frames, showing the $2.05\times$
mean enhancement in the outermost vacuum-facing LLZO layer. [TODO: build from
the existing per-atom error arrays.]

\paragraph{Figure 5.} Pre-round latent-space posterior variance versus the
force error subsequently measured against DFT-FE on the newly labelled
structures ($r=+0.85$), with the selected and random subsets distinguished.

\paragraph{Figure 6.} \label{fig:msd} [TODO --- pending multi-seed runs.]
Mean-squared displacement (in-plane) versus time for the Li-metal and LLZO
regions at 1100~K, with fitted diffusion coefficients and literature
comparison.

% ===========================================================================
\section*{Tables}
% ===========================================================================
% Table~\ref{tab:force-trend} appears inline above per this template's
% layout; house style calls for tables to be submitted as separate,
% editable files at actual submission.

% ===========================================================================
\section*{Figures}
% ===========================================================================
% Figures must be uploaded as separate files at submission (not embedded).

\end{document}
